\section{Ten common claims about ancestral selection, checked}\label{sec:audit}

\noindent A familiar account of how mating shaped the human mind runs as follows. As groups grew, subordinate males became
worth keeping as fighters; men learned to best one another without lasting harm; watching them wrestle, hunt and play revealed
health, drive and intelligence early in life; women and their families chose accordingly; and because the criteria came
before, not after, a lifetime's outcomes, the mind evolved faster. Women in turn were chosen on cues of fertility and health.
Versions of this account circulate widely, and it motivated the present work. It is easy to believe and hard to test, so
each claim is stated and checked in Table~\ref{tab:audit}. Three conclusions shape the rest of the paper.

\textbf{What survives.} The strongest part is the part about hunters (claim N8). Foraging skill is valuable, slow to
mature, variable among men, and hard to rank from a few observations; people in foraging societies say they care about
intelligence, persistence and hard work in a spouse; and better hunters have more children. That is the setting in which
early, repeated, public information should matter most. The claim about subordinate males as fighters (N1) also survives, and
Proposition~\ref{prop:despot} gives it a precise form.

\textbf{What is wrong.} Four things. Restrained fighting is not a human invention and does not require intelligence
(N2): it is the rule among animals. Contests of force select formidability, and comparative data give no sign that stronger
male contest enlarges brains (N3). Formal games are a few thousand years old and cannot have caused a change that was
complete hundreds of thousands of years earlier (N7; Section~\ref{sec:games}). And the cues by which men are said to choose
women (N9, N10) have weak validity as measured and aim largely at a life stage, which is not heritable.

\textbf{What changes in the theory.} The account treats public contest as a single thing that both intensifies
competition and reveals quality. The audit separates four quantities that the theory keeps apart: how accurately a
performance reveals a trait, how strongly standing governs reproduction, which trait the performance reveals, and whether the
performance can be repeated. Leveling within groups lowers the second; safe, varied performance raises the first and the
fourth; and the shift from dominance by force to standing by skill across many domains changes the third. The last of these,
not any increase in competition, is what points selection at a general-purpose brain.

\begin{table*}[t]
\caption{Ten common claims about ancestral selection, checked. Each is stated in the strong form in which it circulates.}
\label{tab:audit}
\scriptsize
\begin{tabularx}{\textwidth}{@{}P{0.022\textwidth}P{0.165\textwidth}P{0.085\textwidth}LP{0.20\textwidth}@{}}
\toprule
 & Claim & Verdict & Evidence and correction & Consequence for the theory\\
\midrule
N1 & As groups grew, subordinate males became valuable as extra fighters, so killing or maiming every rival became costly. &
Supported, restated &
Chimpanzees enter intergroup contests when they outnumber rivals, and killers outnumber victims by a median of eight to
one~\cite{wilson2002,wilson2014}; raids approximate Lanchester's square law~\cite{johnson2015}. Human male reproductive skew is low
for a mammal, in part because men cooperate~\cite{ross2023}, and foragers actively level would-be
dominants~\cite{boehm1993,boehm1999,gavrilets2012}. Corrections: what makes subordinates valuable is the strength of rival coalitions,
not the size of one's own group; ``alpha'' and ``beta'' misdescribe human status, which has a dominance route and a prestige
route~\cite{henrich2001,cheng2013}; and the antiquity and frequency of war are disputed~\cite{bowles2009,fry2013,gomez2016,lahr2016,sala2015}. &
Proposition~\ref{prop:despot}. Leveling lowers intensity (Proposition~\ref{prop:identity}), so by itself it slows selection; its
importance is what it does to the aim.\\
\addlinespace
N2 & Besting rivals without lasting harm became important; boxing or wrestling takes far more intelligence than two rams
butting horns. & Half right &
Restrained, ritualized contest is the animal norm and needs little
cognition~\cite{maynardsmith1973,parker1974,enquist1983,arnott2009}; the rams already do it. What is distinctive in humans: weapons
and coalitions make unregulated fights deadly for either party~\cite{wrangham2018,gintis2015}; contests are rule-bound, refereed,
watched and discussed; and they come in many kinds. Coalitional play fighting is recorded in all 46 forager culture clusters
with adequate descriptions~\cite{scalise2018}, and combat sports are almost exclusively male across 50
societies~\cite{deaner2013}. &
Proposition~\ref{prop:repeat}: a safe contest can be repeated, and repetition is information. The gain is accuracy, not
cleverness in the fighting itself.\\
\addlinespace
N3 & Watching men fight or wrestle revealed health, drive and intelligence, and became an early criterion in mating. &
Strength yes; intelligence weakly &
Observers judge strength and fighting ability accurately from bodies and faces~\cite{sell2009}; Sereer wrestlers have more
children~\cite{llaurens2009}; in one well-studied sample men's dominance, not their attractiveness to women, predicted mating
success~\cite{hill2013}, and contest may have been the main mode of sexual selection in men~\cite{puts2010}. Bystanders use
contests in mate choice in other species, though not always in favor of winners~\cite{doutrelant2000,ophir2003}. But across
primates, stronger male contest goes with body and canine dimorphism~\cite{plavcan2012} and with smaller relative
brains~\cite{schillaci2006}; across mammals it does not enlarge them~\cite{lemaitre2009,fitzpatrick2012}. &
Proposition~\ref{prop:aim}: a contest of force selects formidability. It bears on capacity only through its small loading on
skill.\\
\addlinespace
N4 & Because these were criteria applied before, not at the end of, life, the evolution of mind accelerated. &
Core idea, reworded &
Selection never waits for the end of life. What matters is how accurately those making consequential decisions (marriage,
alliance, support) can forecast a candidate's eventual value at the time they decide. Earlier information raises that
accuracy; in dairy cattle the same move raised genetic gain by 50--100\%~\cite{garciaruiz2016}. ``Accelerated'' needs care:
per generation the change was slow (Section~\ref{sec:explananda}). & Proposition~\ref{prop:identity}.\\
\addlinespace
N5 & Humans selected hero warriors: risk-takers with purpose, strength, and the intelligence to rally others. & Mixed &
Killers had more wives and children among Yanomam\"o~\cite{chagnon1988}, as did prolific raiders among Nyangatom elders, through
bridewealth~\cite{glowacki2015}; but the more aggressive Waorani warriors left fewer descendants~\cite{beckerman2009}. Societies
that offer warriors more rewards lose more men to war~\cite{glowacki2013}. Heroic, not gratuitous, risk is attractive in a
mate~\cite{farthing2005}. Theory predicts that
high-ranking men contribute most to between-group contests~\cite{gavrilets2014}. &
Courage in defense is one performance venue, and standing is how a group pays for a public good. Sign and size depend on
context.\\
\addlinespace
N6 & Women selected tribal leaders for power. & Partly &
Men's status predicts reproductive success across 33 nonindustrial populations~\cite{vonrueden2016}, and women weigh status
more than men do~\cite{buss1989,walter2020}. But forager leaders have little coercive power~\cite{boehm1999}; kin arrange most
forager marriages~\cite{walker2011,apostolou2007}; leaders in the ethnographic record are marked more by expertise and ``neural
capital'' than by dominance~\cite{garfield2019}; and among Tsimane both prestige and dominance pay, prestige chiefly through
wives who begin reproducing earlier~\cite{vonrueden2011}. Reproductive monopoly by powerful men is a Holocene
pattern~\cite{karmin2015,zeng2018} that enlarged no brains. &
The standing route ($\beta_S$), with assessors who are mostly kin. Power detached from assessed capacity is the decoupling
case (prediction~P8).\\
\addlinespace
N7 & Early games displayed the best players in public. & Too recent to be a cause &
Every attested formal game is younger than about 8,000 years, and all secure ones are younger than about 5,500
(Table~\ref{tab:games}): the last 0.2\% of the record, long after brains reached modern size. Games show the propensity. Play,
contest and coalitional play fighting, which are widely documented~\cite{roberts1959,deaner2013,scalise2018,pellis2000}, are the
ancestral analogue. & The mechanism needs public, repeated, outcome-revealing performance, not boards or stadia.\\
\addlinespace
N8 & Successful hunters needed determination, intelligence and health, and both sexes had early indicators of them. &
Supported; the strongest part &
Hunting skill peaks near age~33 and takes ten to twenty years beyond adult size to mature~\cite{koster2020,gurven2006}; better
hunters have more children~\cite{smith2004}. Hadza name ``intelligence'' and ``heart'' beside strong arms as what makes a
hunter~\cite{apicella2014}; Hadza women value intelligence and foraging in a husband, and men value hard work in a
wife~\cite{marlowe2004}. Yet reputations track measured skill only noisily~\cite{stibbardhawkes2018}, and daily returns vary too much
that large samples are needed to tell hunters apart by skill~\cite{hillkintigh2009}. &
The setting of the model: valuable, slow to reveal, noisy. Repetition, diversity and years of observation make it assessable
(Propositions~\ref{prop:portfolio} and~\ref{prop:repeat}; Table~\ref{tab:indicators}).\\
\addlinespace
N9 & Women were chosen for signs that they could bear many children: large hips, large breasts, good health. &
Largely unsupported as stated &
Men's preferences track youth and nubility~\cite{kenrick1992,lassek2019,marlowe1998}. The low waist--hip ratios and body-mass
indices judged most attractive, long read as signs of fecundity~\cite{singh1993}, do not indicate higher fertility or better
health~\cite{lassek2018a,lassek2018b,bovet2019}. A systematic review of direct fertility measures (19 articles, 68 effect
sizes) concluded that the evidence is too weak to support the claim that feminine morphology indicates reproductive potential:
a less feminine waist--hip ratio may reflect past births, and no study has related fertility to facial
femininity~\cite{lidborg2025}. One study
links larger breasts and narrow waists to higher estradiol~\cite{jasienska2004}. Where hips and fat do bear on brains is energy
and obstetrics: infant fat buffers the brain's fuel supply~\cite{kuzawa1998,cunnane2003}, gluteofemoral stores may supply
neurodevelopmental lipids~\cite{lassek2008}, and pelvis, stature and head size covary in ways that ease
birth~\cite{fischer2015,mitteroecker2016,dunsworth2012}. &
Corollary~\ref{cor:cue}: a cue to a life stage is not a heritable target and adds nothing to selection on capacity. The link
to brains is on the budget side ($\kappa_B$).\\
\addlinespace
N10 & Facial symmetry and other facial markers became health indicators. & Preference real; validity weak &
Symmetric and average faces are preferred~\cite{rhodes2006}. The pooled association of fluctuating asymmetry with health and
quality is positive but small, in a literature with publication bias~\cite{vandongen2011}; in 4,732 British adolescents facial
asymmetry was unrelated to childhood health~\cite{pound2014}. The preference may serve to avoid marked developmental disturbance
rather than to rank healthy partners. &
A persuasive cue with low validity. Validity is a parameter, not an assumption (Corollary~\ref{cor:score}).\\
\bottomrule
\end{tabularx}
\end{table*}

